\section{Nonlinear singleton inheritance}
\label{sec:nonlinear-singleton-inheritance}
A surviving singleton contributes an exponential of an old regular
activity. Its first order part must be retained as a Gaussian reference
observable. The error after this subtraction is smaller for two distinct
reasons: the exponential remainder is quadratic, and the actual Gaussian
law is close to the reference law. We prove both estimates on the same
coordinates and then sum them over the inherited occurrences.

\subsection{The nonlinear comparison on fixed coordinates}
Use the notation and quantitative hypotheses of
Section~\ref{sec:principal-gaussian-comparison}. In particular, $R_6,R_9$
are density ratios with the same compensated Jacobian, $\gamma_{C_0}$
is the fixed reference probability measure, and $w(z)$ is the original
masked activity. Thus $w$ is measurable and $|w(z)|\le\tau$. Assume
\[
 0\le\tau\le K\le1.
\]
Write $c(y,z)$ and $v(y,z)$ for the actual and reference cutoffs,
respectively. They take values in $[0,1]$ and have the common compact
support and local plateau properties specified there. The exact integrand
for the singleton plus its reference observable is
\begin{equation}\label{eq:singleton-nonlinear-integrand}
 T(y,z)=R_6(y,z)\{e^{-w(z)c(y,z)}-1\}
                         +R_9(y,z)v(y,z)w(z).
\end{equation}
The plus sign in the second term cancels the negative linear term of
the first. The reference guard ensures that, wherever it is nonzero,
the actual cutoff is one and the frozen field is healthy on the
insertion support. Consequently the same $w$ occurs in both terms,
by the compensated locality identity. No derivative of $w$ is needed.

Let $P_0,P_1,P_2$ denote fixed positive polynomials in the local
coordinate count and jet bounds, enlarged below when necessary. The
density comparison already proved gives
\[
 \|R_6(\cdot,z)\|_{J^n_0}\le P_0,
 \qquad
 \|R_9(\cdot,z)-R_6(\cdot,z)\|_{J^n_0}\le\delta P_1.
\]
Here and below the jet norm includes derivatives of every order from
zero through $n$. The polynomial constants depend on $n$, not on the
number of occurrences or the torus volume. If $Q_c\ge1$ bounds the
cutoff jets, impose also $KQ_c\le1$. For a nonempty derivative set $I$,
the derivative of $e^{-wc}$ is
\[
 e^{-wc}\sum_{\pi\in\mathcal P(I)}
                  \prod_{A\in\pi}(-w D_Ac).
\]
Since $|w|Q_c\le1$, every product has size at most $\tau Q_c$,
up to the fixed number of partitions. For the zero-order term,
$|e^{-wc}-1|\le e\tau$. Product differentiation in
\eqref{eq:singleton-nonlinear-integrand} therefore gives the global
bound $\|T(\cdot,z)\|_{J^n_0}\le\tau P_2$ on the compact read ball.

On the smaller ball where both cutoffs have germ one at $y=0$,
the exact identity is
\[
 T=R_6(e^{-w}-1+w)+(R_9-R_6)w.
\]
For real $|w|\le1$, the exponential series gives
\[
 |e^{-w}-1+w|
 \le |w|^2\sum_{k\ge2}\frac1{k!}
 \le |w|^2\le\tau K.
\]
The two factors involving $w$ are constant in the differentiation
parameter. Hence on this ball
\begin{equation}\label{eq:singleton-good-price}
              \|T(\cdot,z)\|_{J^n_0}
                    \le\tau(KP_0+\delta P_1).
\end{equation}
Let its radius be $v_0$ in the coordinate maximum norm and suppose
$C_0\le MI$. The coordinate Gaussian tail bound used in the principal
comparison gives
\[
 \gamma_{C_0}\{\|z\|_\infty\ge v_0\}
                      \le 2N\exp(-v_0+M/2).
\]
The same compact-coordinate domination justifies differentiating
\eqref{eq:singleton-nonlinear-integrand} under the integral: on a
common parameter neighborhood the smooth factors and their finitely
many derivatives are bounded on the fixed compact ball, and $w$ is
bounded. Splitting that integral into the good ball and its complement
proves
\begin{equation}\label{eq:singleton-comparison-price}
 \left\|\int T(\cdot,z)\,d\gamma_{C_0}(z)\right\|_{J^n_0}
 \le \tau\{KP_0+\delta P_1
                +2N e^{-v_0+M/2}P_2\}.
\end{equation}
This estimate includes the entire near-weight on the bad set; it does
not use cancellation outside the common plateau.

In the renormalization step, $K$ is the incoming regular budget, of
upper exponent $-11/10$. The local covariance and mean comparison of
Section~\ref{sec:principal-comparison-jets} gives $\delta$ of upper
exponent $-1/2$: its two terms are the covariance discrepancy and the
unscaled mean discrepancy multiplied by the inverse noise scale.
The local coordinate counts, cutoff bounds, and polynomials above have
upper exponent zero. The plateau radius is a fixed positive multiple
of the fine monitor and its Gaussian tail has every negative upper
exponent. Thus the braces in \eqref{eq:singleton-comparison-price}
have upper exponent $-1/2$. In particular, all the smallness conditions
in the density comparison and $KQ_c\le1$ hold after increasing the
entrance coupling, with the block size already fixed. Denote a common
nonnegative bound for the braces by $b(x)$.

\subsection{The retained reference term}
We spell out the summation cost of the reference term, since a
per-occurrence gain alone does not give the required coarse budget.
Let $X_i$ be the fine support of occurrence $i$, $a_i$ its fine anchor,
$Y_i$ its allocated coarse support, and $\bar a_i$ the parent of $a_i$.
Put $q_i=\operatorname{cov}_f(X_i)$ and
$q_i'=\operatorname{cov}_c(Y_i)$ for their cover numbers. Suppose the
incoming amplitudes satisfy
\begin{equation}\label{eq:singleton-incoming-budget}
 \sup_a\sum_{i:a_i=a}e^{Aq_i}\tau_i\le E_{\rm in},
                         \qquad A>0.
\end{equation}
All sums use the given occurrence enumeration, including repeated
supports. Fix derivative order $n\ge2$, spatial weight order $p$, and
allocation radius $r_0$. Let $t\ge0$ be the full requested coarse
support exponent, including any reserve needed later.

The reference attenuation proof in
Section~\ref{sec:principal-reference-attenuation} applies to each
forest and tag. In its notation the pointwise derivative bound is
\[
 B_0\tau_i d_L(q_i)
             \{1+\rho_i^{2n}(1+\|z\|)^{4n}\},
 \qquad
 d_L(q)=\begin{cases}L^{-5},&q^3<L,\\1,&q^3\ge L.\end{cases}
\]
Here $B_0$ is fixed before $L$. The same local Gaussian moment
estimate bounds the integral of the polynomial factor. More precisely,
write $H=(r_0+48)(1+B)$, where $B$ is the scale radius. There are fixed
integers $a,v$ such that
\[
 \rho_i\le K_L(\theta_f+g/\theta_f)
                         [H(1+q_i)]^a,
 \quad
 \int(1+\|z\|)^{4n}\,d\gamma\le M_L[H(1+q_i)]^v.
\]
The fine and coarse thresholds here have the bounded ratio required
by that proof. With $h=2na+v$, impose the explicit post-$L$ condition
\[
 [K_L(\theta_f+g/\theta_f)]^{2n}M_L H^h\le1.
\]
It is available by the threshold and monitor argument proved there;
no negative power rate is assigned to $g/\theta_f$. Integration yields
$B_0\tau_i d_L(q_i)\{1+(1+q_i)^h\}$. Since
\[
 d_L(q)\le L^{-5}(1+q)^{15},
\]
including when $q^3\ge L$, the reference price is at most
$2B_0\tau_i L^{-5}(1+q_i)^{15+h}$.

The support compression estimate gives, provided
$20t\,t_{r_0}\le LA$,
\[
 tq_i'\le t\,t_{r_0}+\tfrac12 Aq_i.
\]
Choose $R$ so that $(1+q)^{15+h}\le R e^{Aq/2}$ for $q\ge0$.
Such an $R$ exists by polynomial versus exponential growth and depends
only on the displayed pre-$L$ data. Thus for
$C'=1+2B_0R e^{t t_{r_0}}$, the weighted reference price is bounded by
\begin{equation}\label{eq:singleton-reference-occurrence}
            e^{tq_i'}s_i\le C'L^{-5}e^{Aq_i}\tau_i.
\end{equation}
A finite average over forests and tags whose sum of absolute weights
is at most one preserves this estimate. This is the actual fixed
average used in the singleton construction; no new Gaussian marginal
or smaller-period replacement is introduced.

Each coarse site has exactly $L^4$ fine parents in four dimensions.
For a coarse site $\bar a$, summing
\eqref{eq:singleton-reference-occurrence} first over each fine anchor
above $\bar a$ and then over those $L^4$ anchors proves
\begin{equation}\label{eq:singleton-reference-anchored}
 \sup_{\bar a}\sum_{i:\bar a_i=\bar a}e^{tq_i'}s_i
                                  \le (C'/L)E_{\rm in}.
\end{equation}
The factor $L^4$ is essential. It is paid by the fifth power in the
per-occurrence attenuation.

\subsection{Inheritance of the small remainder}
The corresponding nonlinear remainder is the actual singleton plus
the positive reference observable. Equation
\eqref{eq:singleton-comparison-price} bounds its local ordered jets
by $\tau_i b(x)$ for each fixed forest and tag. The common finite
average preserves this bound.

For completeness, the coordinate-to-field step uses the same local
variation as the actual singleton, not a separately chosen test
field. Restrict each direction to $Y_i$, which leaves every derivative
of a $Y_i$-local function unchanged. For $q_i'\le B$ the weighted
admissibility bounds imply the local direction bound used in the
Gaussian comparison. The support cover restriction and the coarse threshold bound by twice
the fine scale threshold give the required direction norm bound
directly. Indeed the anchor belongs to $Y_i$ and the cover restriction
bounds each support-point distance from it by the local radius $D_x$;
thus the truncated direction satisfies
\[ \|U\|\le\theta_c(1+D_x)^p
                    \le2\theta_f(1+D_x)^p. \]
Apply the comparison to the ordered parameter map formed
from these actual restricted directions. Its mixed derivative is the
prescribed ordered field derivative: evaluate the corresponding
multilinear differential on the coordinate unit vectors, each of norm
one. This conversion incurs no inverse rescaling factor.
Consequently the
estimate controls the prescribed spatially weighted derivatives
through order $n$ on healthy configurations whenever $q_i'\le B$.
Regularity, gauge invariance, and locality on every support follow
from the actual singleton regularity and the reference construction.
The latter is invariant under the orthogonal colour-coordinate gauge
change; locality follows from agreement on $Y_i$. Its measurable
extension outside the healthy domain has the same local germ there.

For the remainder use the weaker consequence
$tq_i'\le t t_{r_0}+Aq_i$ of support compression. Combining it with
\eqref{eq:singleton-incoming-budget} and again counting $L^4$ fine
anchors above a coarse anchor gives
\begin{equation}\label{eq:singleton-remainder-anchored}
 \sup_{\bar a}\sum_{i:\bar a_i=\bar a}
       e^{tq_i'}\tau_i b(x)
       \le L^4 e^{t t_{r_0}}b(x)E_{\rm in}.
\end{equation}
For fixed $L,r_0,t$, the coefficient on the right retains upper
exponent $-1/2$. Equations
\eqref{eq:singleton-reference-anchored} and
\eqref{eq:singleton-remainder-anchored} are the two required
inheritance bounds. The first applies to every reference support;
the derivative assertion for the remainder retains the restriction
$q_i'\le B$. The pre-existing full singleton regularity and the
common funded path are inputs to this step. The argument does not
replace their construction or the estimates for clusters containing
more than one old occurrence.
